\documentclass[a4paper]{jlreq}

\usepackage[haranoaji]{luatexja-preset}
\usepackage[margin=25truemm,headheight=18pt]{geometry}
\usepackage{fancyhdr}
\usepackage{fancyvrb}
\usepackage{graphicx}
\usepackage[section]{placeins}
\usepackage{lastpage}
\usepackage{xcolor}
\usepackage[
  backend=biber,
  style=numeric,
  sorting=none,
  defernumbers=true,
  url=true,
  doi=false,
  isbn=false
]{biblatex}
\addbibresource{bib/github-account.bib}
\usepackage[colorlinks=true,linkcolor=blue!50!black,urlcolor=blue!70!black]{hyperref}

\setlength{\headheight}{18pt}
\pagestyle{fancy}
\fancyhf{}
\fancyhead[L]{GitHub アカウント作成}
\fancyhead[R]{\thepage/\pageref{LastPage}}
\renewcommand{\headrulewidth}{0.4pt}

\DefineVerbatimEnvironment{codeblock}{Verbatim}{
  fontsize=\small,
  frame=single,
  xleftmargin=1.5em,
  xrightmargin=1.5em
}

\newcommand{\cmd}[1]{\texttt{#1}}
\newcommand{\term}[1]{\textbf{#1}}
\newcommand{\materialtitle}{GitHub アカウント作成手順}

\title{\materialtitle}
\author{kska6}
\date{\today}

\begin{document}
\maketitle
\tableofcontents
\newpage

\section*{この資料の目的}

この資料では、GitHub を初めて使う人に向けて、個人アカウントの作成、ログイン後の基本画面の確認、GitHub CLI の認証、学生向け特典を説明します。
GitHub の UI は更新によって変わることがありますが、アカウント作成の基本の流れは「サインアップ画面を開く」「必要な情報を入力する」「メールアドレスを確認する」「ログイン後の画面を確認する」です。

\section{用意するもの}

アカウント作成の事前準備とは、登録に使用するメールアドレス、ユーザー名、認証方法をあらかじめ用意することです。
次のものを用意します。

\begin{itemize}
  \item 受信できるメールアドレス
  \item GitHub で使用するユーザー名
  \item 十分に長いパスワード、または Google / Apple アカウントでのサインアップ
\end{itemize}

ユーザー名は、GitHub 上で自分を識別するための名前です。
リポジトリやプロフィールの URL にも含まれるため、あとから見ても分かりやすい名前にします。
パスワードを使用する場合は、他のサービスと使い回さないようにします。

\section{GitHub のトップページを開く}

ブラウザで GitHub を開きます。
この資料では、主に英語 UI の表記を使用します。

\begin{codeblock}
https://github.com/
\end{codeblock}

画面右上の \term{Sign up}、または画面中央の \term{GitHub に登録する} を選択します。

\begin{figure}[htbp]
  \centering
  \includegraphics[width=0.8\linewidth]{fig/github-home-signup.png}
  \caption{GitHub のトップページとサインアップ入口}
\end{figure}

\section{必要情報を入力する}

アカウント作成とは、サインアップ画面へ必要な情報を入力し、登録を確定する操作です。
メールアドレス、パスワード、ユーザー名、国または地域を入力します。
Google または Apple アカウントを使用する選択肢が表示される場合もあります。

\begin{figure}[htbp]
  \centering
  \includegraphics[width=0.8\linewidth]{fig/github-signup-form.png}
  \caption{GitHub のサインアップフォーム}
\end{figure}

主な入力項目は次のとおりです。

\begin{itemize}
  \item \term{Email}: GitHub からの確認メールを受け取れるアドレスを入力する
  \item \term{Password}: 画面の条件を満たす、安全なパスワードを入力する
  \item \term{Username}: GitHub 上で表示される名前を入力する
  \item \term{Your Country/Region}: 自分の国または地域を選ぶ
  \item \term{GitHub Copilot} や案内メール: 希望しない場合はチェックを外す
\end{itemize}

入力後、\term{Create account} を選択します。
途中で画像認証や追加確認が表示された場合は、画面の指示に従います。

\section{メールアドレスを確認する}

メールアドレスの確認とは、GitHub から届く確認メールの案内に従い、登録したメールアドレスの確認を完了する操作です。
作成したアカウントを利用するため、登録したメールアドレスに確認メールが届いたら、この操作を完了します。

メールが届かない場合は、迷惑メールフォルダ、入力したメールアドレス、受信制限を確認します。
学校や組織のメールアドレスでは、受信制限によって確認メールが届かないことがあります。

\section{ログイン後のダッシュボードを確認する}

ダッシュボードの確認とは、ログイン後の画面で、GitHub の作業を始めるために必要な要素を把握する操作です。
アカウントを作成してログインすると、作業の起点となるダッシュボードが表示されます。
画面には、自分のリポジトリ、ホームやフィード、検索欄、作成メニューなどが表示されます。

最初に見る場所は次の 3 つです。

\begin{itemize}
  \item GitHub ロゴ: ホームに戻る
  \item 検索欄: リポジトリ、ユーザー、Issue、Pull Request を検索する
  \item プロフィール画像: 自分のプロフィールや設定を開く
\end{itemize}

\section{プロフィールを確認する}

右上のプロフィール画像から、自分のプロフィール画面へ移動できます。
プロフィール画面では、ユーザー名、プロフィール画像、公開リポジトリ、活動履歴などを確認できます。

\begin{figure}[htbp]
  \centering
  \includegraphics[width=0.8\linewidth]{fig/github-profile.png}
  \caption{GitHub のプロフィール画面}
\end{figure}

プロフィール画面では、次の点を確認します。

\begin{itemize}
  \item ユーザー名が意図したものになっているか
  \item プロフィール画像や表示名を設定から変更できること
  \item Repositories タブから自分のリポジトリ一覧を見られること
\end{itemize}

\section{作成後にやること}

アカウントを作成したら、次の設定や準備を行います。

\begin{enumerate}
  \item メールアドレスの確認を完了する
  \item プロフィールを必要に応じて整える
  \item 2FA（またはパスキー）を設定してアカウントを保護する
  \item GitHub CLI を使う場合は \cmd{gh auth login} で認証する
  \item 「Git / GitHub 入門」と併用する場合は、リポジトリを作成せず本文資料へ戻る
\end{enumerate}

2FA（二要素認証）は、パスワードに加えて認証アプリやセキュリティキーで本人確認を行う仕組みです。
パスキーは、パスワードに代わる認証方法です。
早めに設定します。

\section{GitHub CLI で認証する}

GitHub CLI は、ターミナルから GitHub を操作するための \cmd{gh} コマンドです。
ブラウザで GitHub にログインしていても、GitHub CLI では別途認証が必要です。
また、HTTPS を使用する Git 操作の認証には、GitHub アカウントのパスワードを使用できません。
GitHub CLI を使用する場合は、次のコマンドで認証します。

\begin{codeblock}
gh auth login
\end{codeblock}

ブラウザ認証とは、ブラウザで一時コードを入力し、GitHub CLI からのアクセスを許可する認証方法です。
選択に迷う場合は、次の設定で進めます。

\begin{itemize}
  \item 接続先は \term{GitHub.com} を選ぶ
  \item Git 操作に使うプロトコルは、迷ったら \term{HTTPS} を選ぶ
  \item 認証方法は、\term{ブラウザ認証}を選ぶ
\end{itemize}

\subsection{表示される選択肢の例}

\cmd{gh auth login} を実行すると、GitHub CLI のバージョンや環境によって文言は異なりますが、次のような質問が表示されます。

\begin{codeblock}
? What account do you want to log into?
> GitHub.com
  GitHub Enterprise Server
\end{codeblock}

通常の GitHub を使う場合は \term{GitHub.com} を選びます。
学校や会社から専用の GitHub Enterprise Server を指定されている場合だけ、そちらを選びます。

\begin{codeblock}
? What is your preferred protocol for Git operations?
> HTTPS
  SSH
\end{codeblock}

初めて設定する場合は、\term{HTTPS} を選びます。
SSH は鍵の作成や登録が必要になるため、授業やチームで指定された場合に選びます。
SSH を選ぶと、既存の SSH 鍵を使用するか、新しい鍵を作成して GitHub に登録するかを確認されることがあります。

\begin{codeblock}
? Authenticate Git with your GitHub credentials?
> Yes
  No
\end{codeblock}

\term{Yes} を選ぶと、GitHub CLI が Git の \cmd{push} や \cmd{clone} に使用する認証情報も設定します。
特別な理由がなければ \term{Yes} を選びます。

\begin{codeblock}
? How would you like to authenticate GitHub CLI?
> Login with a web browser
  Paste an authentication token
\end{codeblock}

初めて認証する場合は、\term{Login with a web browser} を選びます。
\term{Paste an authentication token} は、あらかじめ作成した個人アクセストークンを貼り付ける方法です。
トークンの用途と管理方法を理解している場合に選択します。

\term{Login with a web browser} を選ぶと、次のように一時コードが表示されることがあります。

\begin{codeblock}
! First copy your one-time code: XXXX-XXXX
Press Enter to open github.com in your browser...
\end{codeblock}

表示されたコードをコピーし、Enter を押してブラウザを開きます。
ブラウザの案内に従って GitHub にログインし、コードを入力して GitHub CLI からのアクセスを許可します。
認証が完了したら、次のコマンドで認証状態を確認します。

\begin{codeblock}
gh auth status
\end{codeblock}

認証に成功していれば、\term{github.com} にログインしているアカウントが表示されます。
現在のリポジトリをブラウザで開くには、次のコマンドを使います。

\begin{codeblock}
gh browse
\end{codeblock}

\cmd{gh browse} は、現在のリポジトリに対応する GitHub のページをブラウザで開くコマンドです。
Git リポジトリ内で実行していない場合や、GitHub 上のリポジトリが設定されていない場合は使用できません。

ログアウトしたい場合は、次のコマンドを使います。

\begin{codeblock}
gh auth logout
\end{codeblock}

個人アクセストークンを使用する場合は、ファイルへの保存や画面共有による漏えいに注意します。

\section{学生向け特典（学割）}

学生向け特典（学割）とは、在学状況を証明して申請し、GitHub Pro や \term{GitHub Student Developer Pack} などを利用できるようにする制度です。
GitHub アカウントを作成したあとに、\term{GitHub Education} から申請できます。
特典の内容は変更されることがあるため、申請時に公式ページで最新情報を確認します。

主な申請条件は次のとおりです。

\begin{itemize}
  \item 学位または卒業資格を得られる学校・教育課程に在籍している
  \item 自分の GitHub 個人アカウントを持っている
  \item 13 歳以上である
  \item 学校発行のメールアドレス、学生証、履修表、在学証明書などで在学状況を証明できる
\end{itemize}

申請は、GitHub にログインした状態で次のページから行います。

\begin{codeblock}
https://github.com/settings/education/benefits
\end{codeblock}

画面の \term{Start an application} から必要事項を入力し、在学状況を確認できる資料を提出します。
学校発行のメールアドレスを求められた場合は、そのアドレスを GitHub アカウントへ追加し、確認を済ませてから申請します。

\section{公式資料}

\nocite{*}

\subsection{GitHub アカウントの作成}

\printbibliography[heading=none,keyword=account]

\subsection{学生向け特典}

\printbibliography[heading=none,keyword=education]

\subsection{GitHub CLI の認証}

\printbibliography[heading=none,keyword=cli]

\end{document}
